\chapter{Conclusions}

This project set out to create a tool that would use restrictions to provide a new, effective method for novice programmers to learn Python. Overall, it was a general success. All the core features of the program were implemented, alongside the partial implementations to readable exceptions. Whilst the extension to readable exceptions and the program analyser were abandoned, this was largely because they were too ambitious in scope, and the tool still feels complete without them.

In evaluating the project, it was intended to show that the tool was both effective and usable. The study, while underpowered due to the limitations of the project resources, hints towards promising results should the ideas presented be extended and studied further, and suggests that the tool was more usable than the baseline. Additionally, the timing analysis showed that adding restrictions could have insignificant overhead.

Having completed the project, there were two key takeaways. First, the initial plans underestimated how long some components (such as the visualiser) would have taken to implement, whilst being overly generous with the `catch-up' time. These evened out, but missing ambitious deadlines caused unnecessary stress. Were the project to be completed again, more allowance would have been given to allocated time blocks over free weeks for catching up. Second, more time would have been given to the preparation stage. In the proposal, two weeks were dedicated to this, but three were used. This extra time on preparation made the implementation stage go much more smoothly, but it was underappreciated in the initial plans.

There is clear room to extend the ideas proposed in this project. As mentioned in the preparation chapter, the idea could be extended to block-based languages by restricting to subsets of the available blocks; extending to other text-based languages may also hold value. Furthermore, there are several clear avenues for expansion of the tool in its current form. This project focused on beginner programmers for a clearer target audience and appropriate Part II project scope, but extending the concept to more advanced Python concepts, such as functional or object-oriented programming, could extend the `lifetime' of the tool (i.e. how long learners find it useful for). Alternatively, extensions may attempt to make the tool more granular by adding more options, such as the ability to restrict comprehension separately to lists. While the system could be deployed in its current state, a login system might be useful if the program were used more widely, such as in a school or as part of a course.

